\section{Introduction}
\label{sec:introduction}

Reinforcement learning (RL) has emerged as a prominent paradigm for sequential decision-making in complex, uncertain environments, achieving strong performance across domains ranging from games to continuous control in robotics \citep{mnih2015humanlevel,sutton2018reinforcement}. Financial trading is a particularly compelling yet inherently difficult RL application: decisions are sequential, delayed consequences are material, and market microstructure imposes strict feasibility constraints~\citep{moody2001learning,zhang2020deep,gould2013limit}. Within global financial markets, foreign exchange (Forex) is especially demanding. Unlike many equity settings, Forex trading is continuous (24/5), highly leveraged, and sensitive to financing and transaction frictions, while return distributions remain heavy-tailed and regime-dependent \citep{cont2001empirical}.

Despite substantial recent activity and promising results in applied RL-for-trading research—such as predictive signal extraction \citep{deng2016deep}, multi-asset portfolio allocation \citep{jiang2017deep}, and raw return optimization algorithms \citep{meng2019reinforcement,theate2021application}—three structural gaps continue to limit real-world applicability. First, \emph{environment fidelity} is often insufficient. Real-world financial environments are heavily driven by execution mechanics that are challenging to simulate, yet crucial for performance \citep{zhang2020deep}. Many studies rely on simplified simulators that under-model or omit key factors such as spread-embedded execution friction, dynamic financing costs, margin-based liquidation, and realistic decision-execution timing \citep{carapuco2018reinforcement,rundo2019deep}. These simplifications introduce a significant sim-to-real gap, often yielding strategies that are economically brittle outside controlled simulations.

Second, \emph{reward design opacity} obscures core learning dynamics. Conventional approaches typically optimize a single scalar reward (e.g., Profit and Loss (PnL), equity delta, or risk-adjusted metrics like proxies for the Sharpe ratio) \citep{moody2001learning,sharpe1994sharpe,theate2021application}. While effective for driving generalized learning, such monolithic formulations complicate credit assignment and make it difficult to disentangle which signals improve trade timing versus which penalties suppress excessive trading activity. Without systematic decomposition and ablation, reward shaping remains difficult to functionally analyze, tune, and reproduce \citep{ng1999policy}.

Third, \emph{action-space granularity} is frequently misspecified. Existing formulations often rely on either coarse discrete spaces (e.g., simplistic buy/sell/hold commands) \citep{theate2021application} or unconstrained continuous position sizing derived from deterministic policy gradients \citep{jiang2017deep,schulman2017proximal,haarnoja2018soft}. Coarse spaces fail to capture practical trading operations such as scaling into positions (pyramiding), reducing exposure, or executing explicit reversals. Conversely, continuous formulations can bypass dynamically enforceable constraints such as margin limits and position feasibility, which are often more naturally handled using action masking in structured discrete domains \citep{huang2020closer}.

To address these challenges, this paper presents a unified, open-source framework designed around these three dimensions. The proposed system integrates: (i) a high-fidelity, Gymnasium-compatible Forex environment with strict anti-lookahead execution, realistic friction and rollover modeling, and margin-aware risk controls; (ii) an 11-component decomposable reward architecture with fixed, pre-specified aggregation, state-dependent clipping, and transparent per-step component logging; and (iii) a legal-action-masked, 10-action discrete space with explicit primitives for pyramiding, martingale-style scaling, reduction, closure, and reversal. All experiments are driven through configuration-based protocols, enabling controlled generation of experimental variants without ad-hoc code modifications.

\paragraph{Contributions.}
Our primary contributions are methodological and infrastructural:
\begin{itemize}[leftmargin=1.4em]
    \item \textbf{Decomposable reward architecture:} We introduce an 11-component reward system for Forex RL, where distinct economic drivers and penalties are independently enabled, weighted, and logged, supporting systematic ablation and component-wise attribution analysis.
    \item \textbf{Structured action space with legality masking:} We formulate a 10-action discrete interface governed by margin-aware legality constraints, applying valid-action masking during both environment interaction and replay-based learning.
    \item \textbf{High-fidelity execution environment:} We design an environment with explicit anti-lookahead timing (observe close$_t$, execute open$_{t+1}$, mark close$_{t+1}$), comprehensive transaction cost modeling, and realistic rollover financing mechanisms, including triple-swap Wednesdays.
    \item \textbf{Controlled experimental protocol:} We demonstrate the framework through three experiment families (reward ablation, action-space comparison, and scaling analysis) on the EURUSD training split, with all variants generated via configuration rather than structural code changes.
    \item \textbf{Reproducible research infrastructure:} We provide a complete pipeline enabling deterministic reproducibility under controlled conditions through resolved configuration snapshots, fixed seeding, isolated artifact management, and leakage-free feature processing.
\end{itemize}

The empirical emphasis of this work is methodological robustness rather than state-of-the-art performance benchmarking. Reproducibility and experimental scope (including seeding and evaluation partitioning) are detailed in Section~\ref{sec:experiments}; interpretations focus on system behaviour and learning dynamics rather than stochastic performance variation, consistent with cautionary guidance in trading-RL evaluation~\citep{zhang2020deep}.

\paragraph{Research questions.}
The empirical program is organised around three focused research questions:
\begin{enumerate}[leftmargin=1.8em, label=\textbf{RQ\arabic*.}]
    \item Which reward components contribute useful training signal, and does progressive penalty accumulation produce non-monotonic effects on training-period risk-adjusted return?
    \item Does increased action-space granularity improve training efficiency and risk-adjusted learning relative to a coarse three-action discrete formulation?
    \item How do asymmetric position-scaling strategies (pyramiding vs.\ martingale-style scaling) influence training-period risk profiles, drawdown dynamics, and average scaling behaviour?
\end{enumerate}

See Section~\ref{sec:experiments} for the canonical experimental setup, including evaluation partitioning and reproducibility controls.

The remainder of the paper is organized as follows. \Cref{sec:related_work} situates the framework within existing RL trading literature and identifies the research gap. \Cref{sec:methodology} details the environment, reward, and agent design. \Cref{sec:experiments} describes the evaluation protocol and reproducibility controls. \Cref{sec:results} presents the empirical findings, and \Cref{sec:discussion} discusses implications. \Cref{sec:future_work} then outlines limitations and forward directions before \Cref{sec:conclusion} concludes.